\chapter[33]{Noncompact Hyperbolic Limits}
\label{ch:chow-iv-noncompact-hyperbolic-limits}
\dcref{ch33}
\source{chow-et-al-ricci-flow-part-iv:page-0234}

Throughout this chapter, \((\mathcal M^3,g(t))\) is a connected, closed,
orientable nonsingular solution of the normalized Ricci flow in Case III which
has a complete noncompact hyperbolic limit.  A finite-volume hyperbolic limit
\((\mathcal H^3,h)\) is connected, orientable, complete, and noncompact.  For
\(A>0\), write \(\mathcal H_A\) for the compact truncation obtained by cutting
each standard cusp at its torus of area \(A\).

\section{Hyperbolic pieces}

\begin{definition}[CMC boundary conditions]
\label{def:chow-iv-cmc-boundary-conditions}
\source{chow-et-al-ricci-flow-part-iv:page-0236}
\dcref{ch33:33.1}
\group{Hyperbolic pieces}

Let \(F:(\mathcal H_A,h)\to(\mathcal M,g)\) be an embedding.  It satisfies the
\emph{CMC boundary conditions} when every component of
\(F(\partial\mathcal H_A)\) has constant mean curvature and area \(A\), and
\(F_*N\) is normal to \(F(\partial\mathcal H_A)\), where \(N\) is the outward
unit normal of \(\partial\mathcal H_A\) with respect to \(h\).
\end{definition}

\begin{definition}[Stable hyperbolic limit]
\label{def:chow-iv-stable-hyperbolic-limit}
\source{chow-et-al-ricci-flow-part-iv:page-0237}
\uses{def:chow-iv-cmc-boundary-conditions}
\dcref{ch33:33.2}
\group{Hyperbolic pieces}


A finite-volume hyperbolic limit \((\mathcal H^3,h)\) is \emph{stable} if there
are a time \(T_0\), a smooth nonincreasing function
\(A:[T_0,\infty)\to(0,\widetilde A]\) tending to zero, compact submanifolds
\(\mathcal M_{A(t)}(t)\subset\mathcal M\) with torus boundary, and harmonic
diffeomorphisms
\[
 F(t):(\mathcal H_{A(t)},h)\longrightarrow(\mathcal M_{A(t)}(t),g(t))
\]
satisfying the CMC boundary conditions, such that for every \(m\in\mathbb N\),
\[
 \lim_{t\to\infty}\|F(t)^*g(t)-h\|_{C^m(\mathcal H_{A(t)},h)}=0.
\]
The images form a \emph{stable asymptotically hyperbolic submanifold}.
\end{definition}

\begin{remark}[Properties of stable hyperbolic limits]
\label{rem:chow-iv-stable-hyperbolic-limit-properties}
\source{chow-et-al-ricci-flow-part-iv:page-0237}
\uses{def:chow-iv-stable-hyperbolic-limit}
\dcref{ch33:33.3}
\group{Hyperbolic pieces}


The boundary and the side selected by \(\mathcal M_{A(t)}(t)\) vary smoothly;
the pieces converge pointedly and smoothly to \((\mathcal H,h)\); fixed
truncation tori become concave and asymptotically totally umbilic.  The stable
submanifold itself need not be unique.
\end{remark}

\begin{definition}[Immortal almost hyperbolic piece]
\label{def:chow-iv-immortal-almost-hyperbolic-piece}
\source{chow-et-al-ricci-flow-part-iv:page-0237,chow-et-al-ricci-flow-part-iv:page-0238}
\uses{def:chow-iv-cmc-boundary-conditions}
\dcref{ch33:33.4}
\group{Hyperbolic pieces}


For \(A\in(0,\sqrt3/4]\) and \(k\in\mathbb N\), an \emph{immortal
\((A,k)\)-almost hyperbolic piece} is a smooth family
\(\mathcal M_A(t)\), \(t\ge T\), together with a finite-volume hyperbolic
manifold \((\mathcal H,h)\) and harmonic diffeomorphisms
\(F(t):\mathcal H_A\to\mathcal M_A(t)\) satisfying the CMC boundary conditions
and
\[
 \|F(t)^*g(t)-h\|_{C^k(\mathcal H_A,h)}\le k^{-1}
 \quad(t\ge T).
\]
\end{definition}

\begin{proposition}[Immortality of almost hyperbolic pieces]
\label{prop:chow-iv-almost-hyperbolic-pieces-immortal}
\source{chow-et-al-ricci-flow-part-iv:page-0238}
\uses{def:chow-iv-immortal-almost-hyperbolic-piece, prop:chow-iv-harmonic-parametrizations, prop:chow-iv-continue-harmonic-parametrizations}
\dcref{ch33:33.5}
\group{Stability of hyperbolic limits}


If pointed time translates based at \((x_i,t_i)\), with \(t_i\to\infty\),
converge to a finite-volume hyperbolic limit \((\mathcal H,h,x_\infty)\), then
for every admissible \(A\) and every \(k\), some sufficiently late almost
hyperbolic piece extends to an immortal \((A,k)\)-almost hyperbolic piece
corresponding to \((\mathcal H,h)\).
\end{proposition}

\begin{proposition}[Stability of hyperbolic limits]
\label{prop:chow-iv-stability-hyperbolic-limits}
\source{chow-et-al-ricci-flow-part-iv:page-0238}
\uses{def:chow-iv-stable-hyperbolic-limit, prop:chow-iv-almost-hyperbolic-pieces-immortal, thm:chow-iv-minimal-cusp-limit-stable, lem:chow-iv-almost-hyperbolic-pieces-disjoint}
\dcref{ch33:33.6}
\group{Stability of hyperbolic limits}


Every finite-volume hyperbolic limit of the nonsingular solution is stable.
\end{proposition}

\begin{definition}
\label{def:chow-iv-immortal-asymptotically-hyperbolic-piece}
\source{chow-et-al-ricci-flow-part-iv:page-0238,chow-et-al-ricci-flow-part-iv:page-0239}
\uses{def:chow-iv-immortal-almost-hyperbolic-piece}
\dcref{ch33:33.7}
\group{Hyperbolic pieces}


An \emph{immortal asymptotically hyperbolic piece} corresponding to
\((\mathcal H,h)\) is a smooth family \(\mathcal M_A(t)\), defined for all
large \(t\), whose boundary consists of CMC tori of area \(A\), and for which
\((\mathcal M_A(t),g(t))\) converges smoothly to \((\mathcal H_A,h)\).
Consequently it is eventually an immortal \((A,k)\)-almost hyperbolic piece
for every fixed \(k\).
\end{definition}

\begin{proposition}[Decomposition into hyperbolic and collapsed pieces]
\label{prop:chow-iv-hyperbolic-collapsed-decomposition}
\source{chow-et-al-ricci-flow-part-iv:page-0239,chow-et-al-ricci-flow-part-iv:page-0254,chow-et-al-ricci-flow-part-iv:page-0255,chow-et-al-ricci-flow-part-iv:page-0256,chow-et-al-ricci-flow-part-iv:page-0257,chow-et-al-ricci-flow-part-iv:page-0258}
\uses{prop:chow-iv-stability-hyperbolic-limits, lem:chow-iv-almost-hyperbolic-pieces-disjoint}
\dcref{ch33:33.8}
\group{Stability of hyperbolic limits}


There is a finite family of complete finite-volume hyperbolic manifolds
\((\mathcal H_\alpha,h_\alpha)\) such that, for every sufficiently small
\(\varepsilon>0\), sufficiently late time slices contain pairwise disjoint
harmonically parametrized truncations of these manifolds, \(\varepsilon\)-close
in \(C^{[1/\varepsilon]}\) and satisfying the CMC boundary conditions, while
the complement has injectivity radius everywhere less than \(\varepsilon\).
\end{proposition}

\begin{proposition}[Incompressibility of boundary tori]
\label{prop:chow-iv-hyperbolic-boundary-incompressible-summary}
\source{chow-et-al-ricci-flow-part-iv:page-0239}
\uses{thm:chow-iv-immortal-piece-boundary-incompressible}
\dcref{ch33:33.9}
\group{Incompressibility}


Every boundary torus of an immortal asymptotically hyperbolic piece is
incompressible in \(\mathcal M\).
\end{proposition}

\begin{definition}[Sweep-out]
\label{def:chow-iv-hypersurface-sweep-out}
\source{chow-et-al-ricci-flow-part-iv:page-0240}
\dcref{ch33:33.10}
\group{CMC and harmonic parametrizations}

A family of hypersurfaces \(\{S_s\}_{s\in I}\) \emph{sweeps out} a region
\(U\) when \(U\subset\bigcup_{s\in I}S_s\).
\end{definition}

\begin{proposition}[Existence of a CMC sweep-out in almost hyperbolic cusps]
\label{prop:chow-iv-cmc-sweepout-almost-cusp}
\source{chow-et-al-ricci-flow-part-iv:page-0240}
\uses{def:chow-iv-hypersurface-sweep-out}
\dcref{ch33:33.11}
\group{CMC and harmonic parametrizations}


Given \([a,b]\subset\mathbb R\) and \(\varepsilon_0>0\), every smooth metric
on \([a,b]\times\mathcal V\) sufficiently close in \(C^{2,\alpha}\) to
\(dr^2+e^{-2r}g_{\mathrm{flat}}\) admits a smooth family of CMC hypersurfaces
which sweeps out \((a+\varepsilon_0,b-\varepsilon_0)\times\mathcal V\) and is
\(C^{2,\alpha}\)-close to the standard slices.
\end{proposition}

\begin{proposition}[Existence of harmonic maps near the identity]
\label{prop:chow-iv-harmonic-maps-near-identity}
\source{chow-et-al-ricci-flow-part-iv:page-0240}
\dcref{ch33:33.12}
\group{CMC and harmonic parametrizations}

Let \((N^n,g)\) be compact with negative Ricci curvature and concave boundary.
For every smooth metric \(\widetilde g\) sufficiently close to \(g\) in
\(C^{2,\alpha}\), there is a unique smooth harmonic diffeomorphism
\(F:(N,g)\to(N,\widetilde g)\), close to the identity, which maps the outward
normal to a \(\widetilde g\)-normal along the boundary.  For smoothly varying
metrics, the resulting maps vary smoothly.
\end{proposition}

\begin{proposition}[Existence of harmonic parametrizations of almost hyperbolic pieces]
\label{prop:chow-iv-harmonic-parametrizations}
\source{chow-et-al-ricci-flow-part-iv:page-0241,chow-et-al-ricci-flow-part-iv:page-0242}
\uses{def:chow-iv-cmc-boundary-conditions, prop:chow-iv-cmc-sweepout-almost-cusp, prop:chow-iv-harmonic-maps-near-identity}
\dcref{ch33:33.13}
\group{CMC and harmonic parametrizations}


Suppose pointed complete Riemannian three-manifolds
\((\mathcal M_i,g_i,x_i)\) converge smoothly in the pointed Cheeger--Gromov
sense to a noncompact finite-volume hyperbolic manifold
\((\mathcal H,h,x_\infty)\).  If \(A\in(0,\sqrt3/8]\) and
\(x_\infty\in\operatorname{int}(\mathcal H_A)\), then for all large \(i\)
there is a harmonic embedding \(F_i:\mathcal H_A\to\mathcal M_i\) satisfying
the CMC boundary conditions, with
\[
 \|F_i^*g_i-h\|_{C^k(\mathcal H_A,h)}\longrightarrow0,
 \qquad d_{g_i}(F_i(x_\infty),x_i)\longrightarrow0
\]
for every \(k\).
\end{proposition}

\begin{proof}
\sketch
Cheeger--Gromov convergence supplies almost-isometric embeddings of the larger
truncation \(\mathcal H_{A/2}\).  The CMC sweep-out produces, in every end, a
CMC torus of area \(A\); these tori bound a domain converging smoothly to
\(\mathcal H_A\).  Pull the metric on that domain back to \(\mathcal H_A\) and
apply the near-identity harmonic-map proposition.  Composing the resulting
harmonic diffeomorphism with the original convergence map gives \(F_i\).
Smooth dependence in the implicit function theorem gives convergence in every
\(C^k\)-norm, and convergence of the auxiliary diffeomorphisms to the identity
gives the assertion about the basepoints.
\end{proof}

\begin{proposition}[Continuing almost-isometric harmonic parametrizations]
\label{prop:chow-iv-continue-harmonic-parametrizations}
\source{chow-et-al-ricci-flow-part-iv:page-0243,chow-et-al-ricci-flow-part-iv:page-0244,chow-et-al-ricci-flow-part-iv:page-0245,chow-et-al-ricci-flow-part-iv:page-0246,chow-et-al-ricci-flow-part-iv:page-0247}
\uses{prop:chow-iv-cmc-sweepout-almost-cusp, prop:chow-iv-harmonic-maps-near-identity, prop:chow-iv-harmonic-parametrizations}
\dcref{ch33:33.14}
\group{CMC and harmonic parametrizations}


For every finite-volume hyperbolic \((\mathcal H,h)\) and admissible \(A\),
there is \(k_0\ge4\) such that an initially harmonic CMC embedding
\(F_\alpha:\mathcal H_A\to(\mathcal M,g(\alpha))\) satisfying
\(\|F_\alpha^*g(\alpha)-h\|_{C^k}\le k^{-1}\), \(k\ge k_0\), extends uniquely
to a smooth family on a maximal interval \([\alpha,\beta]\).  The same bound
holds throughout; either \(\beta=\omega\), or equality holds in the bound at
\(\beta\).
\end{proposition}

\begin{proof}
\sketch
Openness follows by continuing the CMC boundary tori with the sweep-out theorem
and then applying the near-identity harmonic-map theorem to their bounded
domains.  For closedness, the metric bound controls \(dF\) and its inverse.
Differentiating \(F^*g\) expresses each higher covariant derivative of \(F\) in
terms of derivatives of \(F^*g\), lower derivatives of \(F\), and the ambient
metric.  Arzela--Ascoli and elliptic regularity therefore give smooth
subconvergence at an endpoint.  The limit remains an embedding, satisfies the
CMC conditions and the same estimate, and local uniqueness identifies every
continuation.  Thus continuation stops only at the stated threshold.
\end{proof}

\section{Stability of hyperbolic limits}

\begin{definition}[Space of hyperbolic limits]
\label{def:chow-iv-space-hyperbolic-limits}
\source{chow-et-al-ricci-flow-part-iv:page-0247,chow-et-al-ricci-flow-part-iv:page-0248}
\dcref{ch33:33.15}
\group{Stability of hyperbolic limits}

The space \(\mathfrak H(\mathcal M,g(t))\) consists of pointed complete
noncompact hyperbolic three-manifolds arising as smooth pointed limits of time
translates of \((\mathcal M,g(t))\).  Its elements have uniformly bounded
volume and hence a uniformly bounded number of cusp ends.
\end{definition}

\begin{proposition}[Minimal-cusp almost hyperbolic pieces are immortal]
\label{prop:chow-iv-minimal-cusp-pieces-immortal}
\source{chow-et-al-ricci-flow-part-iv:page-0248,chow-et-al-ricci-flow-part-iv:page-0249,chow-et-al-ricci-flow-part-iv:page-0250,chow-et-al-ricci-flow-part-iv:page-0251}
\uses{def:chow-iv-space-hyperbolic-limits, prop:chow-iv-continue-harmonic-parametrizations, cor:chow-iv-mostow-equivalences, prop:chow-iv-negative-case-hyperbolic-limit}
\dcref{ch33:33.16}
\group{Stability of hyperbolic limits}


If \((\mathcal H,h)\in\mathfrak H(\mathcal M,g(t))\) has the minimum number
of cusp ends among all hyperbolic limits, then every sufficiently accurate
almost hyperbolic truncation of \(\mathcal H\) at a sufficiently late time
extends to an immortal almost hyperbolic piece.
\end{proposition}

\begin{proof}
\sketch
Otherwise, choose terminal times at which the continuation estimate first
becomes an equality.  Compactness of pointed normalized Ricci flows produces a
new finite-volume hyperbolic limit.  The terminal parametrizations converge
smoothly to a harmonic almost isometry from \(\mathcal H_A\) into that limit.
Minimality of the number of cusps and Mostow rigidity extend it to an isometry;
uniqueness of the near-identity harmonic map then says that the limiting map is
that isometry.  This contradicts equality at the continuation threshold.
\end{proof}

\begin{theorem}[Stability of hyperbolic limits with minimal cusp ends]
\label{thm:chow-iv-minimal-cusp-limit-stable}
\source{chow-et-al-ricci-flow-part-iv:page-0251,chow-et-al-ricci-flow-part-iv:page-0252,chow-et-al-ricci-flow-part-iv:page-0253,chow-et-al-ricci-flow-part-iv:page-0254}
\uses{def:chow-iv-stable-hyperbolic-limit, prop:chow-iv-minimal-cusp-pieces-immortal, lem:chow-iv-small-flat-torus-in-cusp}
\dcref{ch33:33.17}
\group{Stability of hyperbolic limits}


A hyperbolic limit having the minimum number of cusp ends among all elements of
\(\mathfrak H(\mathcal M,g(t))\) is stable.
\end{theorem}

\begin{proof}
\sketch
Apply immortality at truncation levels \(1/k\).  A volume-packing argument
selects a subsequence whose pieces intersect at later starting times.  The
small-flat-torus lemma forces the coarser truncations to be nested inside the
finer ones, and this nesting persists because their CMC boundaries vary
continuously.  Interpolate the truncation areas by a smooth decreasing function
\(A(t)\to0\).  Near-identity uniqueness then glues the corresponding harmonic
parametrizations smoothly across the selected times, with error tending to zero
in every fixed \(C^m\)-norm.
\end{proof}

\begin{lemma}[Almost-flat small tori lie in cusp regions]
\label{lem:chow-iv-small-flat-torus-in-cusp}
\source{chow-et-al-ricci-flow-part-iv:page-0252,chow-et-al-ricci-flow-part-iv:page-0253}
\uses{cor:chow-iv-small-almost-flat-tori}
\dcref{ch33:33.18}
\group{Stability of hyperbolic limits}


An embedded two-torus in an almost hyperbolic piece, with bounded intrinsic and
extrinsic curvature and sufficiently small area, is contained in an almost
hyperbolic cusp region.
\end{lemma}

\begin{proof}
\sketch
An almost isometry carries the torus to one with the same uniform bounds and
small area in the limiting finite-volume hyperbolic manifold.  The
corresponding hyperbolic cusp result places it in a maximal cusp, whose inverse
image is the required almost cusp region.
\end{proof}

\begin{lemma}[Disjointness of almost hyperbolic pieces]
\label{lem:chow-iv-almost-hyperbolic-pieces-disjoint}
\source{chow-et-al-ricci-flow-part-iv:page-0254}
\uses{lem:chow-iv-small-flat-torus-in-cusp, thm:chow-iv-finite-volume-ends-cusps}
\dcref{ch33:33.19}
\group{Stability of hyperbolic limits}


For sufficiently small truncation areas and sufficiently accurate
parametrizations, two continuously varying almost hyperbolic pieces are either
disjoint, or each doubled truncation is contained in the other piece, throughout
their common interval of existence.
\end{lemma}

\begin{proof}
\sketch
Every sufficiently small boundary torus of either piece lies in a cusp region
of the other whenever the pieces meet.  The standard product structure of
finite-volume hyperbolic ends orders the corresponding truncation tori.  The
ordering gives both inclusions, and continuity prevents it from changing
without a boundary contact, where the same cusp argument applies again.
\end{proof}

\section{Incompressibility of boundary tori}

\begin{theorem}[Incompressibility of boundary tori]
\label{thm:chow-iv-immortal-piece-boundary-incompressible}
\source{chow-et-al-ricci-flow-part-iv:page-0259}
\uses{def:chow-iv-disk-class-almost-piece, prop:chow-iv-compressible-torus-area-decay, lem:chow-iv-boundary-torus-incompressible, lem:chow-iv-separating-incompressible-surfaces}
\dcref{ch33:33.21}
\group{Incompressibility}


For every admissible truncation area and every time at which the corresponding
immortal asymptotically hyperbolic piece is defined, its boundary is
incompressible in \(\mathcal M\).
\end{theorem}

\begin{proof}
\sketch
It is incompressible on the hyperbolic side.  If it were compressible on the
complementary side, the least area \(A_{\mathfrak T_A}(t)\) below would be
positive, while for every \(\varepsilon\in(0,2\pi)\) its upper right derivative
would eventually be at most \(-2\pi+\varepsilon\).  Integration would make
that positive quantity negative in finite time, a contradiction.  The
separating-surface criterion then gives incompressibility in all of
\(\mathcal M\).
\end{proof}

\begin{definition}[Disk class and least area]
\label{def:chow-iv-disk-class-almost-piece}
\source{chow-et-al-ricci-flow-part-iv:page-0259}
\dcref{ch33:33.22}
\group{Minimal disks}

Given a boundary torus of \(\mathcal M_A(t)\), let
\(\mathfrak T_A(t)\) be the smooth maps \(f:D^2\to\mathcal M\setminus
\mathcal M_A(t)\) whose boundary lies in that torus and represents a nontrivial
element of its fundamental group.  Define
\[
 A_{\mathfrak T_A}(t)=\inf_{f\in\mathfrak T_A(t)}
 \operatorname{Area}_{g(t)}f(D^2).
\]
\end{definition}

\begin{remark}[Notation for embedded minimizing disks]
\label{rem:chow-iv-disk-class-notation}
\source{chow-et-al-ricci-flow-part-iv:page-0259}
\uses{def:chow-iv-disk-class-almost-piece}
\dcref{ch33:33.23}
\group{Minimal disks}


The truncation parameter is suppressed when fixed, and for an embedding
\(f\in\mathfrak T_A(t)\) its image disk may itself be denoted as an element of
\(\mathfrak T_A(t)\).
\end{remark}

\begin{proposition}[Compressible tori yield linearly decreasing disk areas]
\label{prop:chow-iv-compressible-torus-area-decay}
\source{chow-et-al-ricci-flow-part-iv:page-0259,chow-et-al-ricci-flow-part-iv:page-0260,chow-et-al-ricci-flow-part-iv:page-0261,chow-et-al-ricci-flow-part-iv:page-0262,chow-et-al-ricci-flow-part-iv:page-0263,chow-et-al-ricci-flow-part-iv:page-0264,chow-et-al-ricci-flow-part-iv:page-0265,chow-et-al-ricci-flow-part-iv:page-0266,chow-et-al-ricci-flow-part-iv:page-0267,chow-et-al-ricci-flow-part-iv:page-0268,chow-et-al-ricci-flow-part-iv:page-0269,chow-et-al-ricci-flow-part-iv:page-0270,chow-et-al-ricci-flow-part-iv:page-0271,chow-et-al-ricci-flow-part-iv:page-0272,chow-et-al-ricci-flow-part-iv:page-0273,chow-et-al-ricci-flow-part-iv:page-0274,chow-et-al-ricci-flow-part-iv:page-0275}
\uses{def:chow-iv-disk-class-almost-piece, thm:chow-iv-meeks-yau, lem:chow-iv-minimizing-disk-evolution, lem:chow-iv-area-decay-short-boundary, prop:chow-iv-minimizing-disk-boundary-small, lem:chow-iv-area-decay-long-boundary}
\dcref{ch33:33.24}
\group{Minimal disks}


If a boundary torus of an immortal asymptotically hyperbolic piece is
compressible in the complementary side, then \(\mathfrak T_A(t)\ne\varnothing\),
\(A_{\mathfrak T_A}(t)>0\), and for every \(\varepsilon>0\), at all sufficiently
large times,
\[
 \frac{d^+A_{\mathfrak T_A}}{dt}(t)
 :=\limsup_{s\to0^+}\frac{A_{\mathfrak T_A}(t+s)-A_{\mathfrak T_A}(t)}s
 \le-2\pi+\varepsilon.
\]
\end{proposition}

\begin{proof}
\sketch
The Meeks--Yau theorem supplies a least-area embedded disk normal to the torus.
If its boundary length is at most the prescribed multiple of its area,
Lemma~\ref{lem:chow-iv-area-decay-short-boundary} gives the result.  In the
complementary case, the long-cusp comparison argument makes the boundary
arbitrarily short, and the minimizing-disk evolution inequality then gives the
same derivative bound.  These two cases exhaust all possibilities.
\end{proof}

\begin{theorem}[Meeks--Yau least-area disk]
\label{thm:chow-iv-meeks-yau}
\source{chow-et-al-ricci-flow-part-iv:page-0260}
\dcref{ch33:33.25}
\group{Minimal disks}

Let \((N^3,g)\) be compact with convex compressible boundary.  Among smooth
disk maps whose boundary represents a nontrivial element of
\(\pi_1(\partial N)\), there is a least-area conformal harmonic map \(f_0\).
It is embedded, normal to the boundary, and its boundary class is primitive.
If that boundary lies in a torus \(T\), it generates the infinite cyclic group
\(\ker(\pi_1(T)\to\pi_1(N))\).
\end{theorem}

\begin{remark}[Relative Meeks--Yau theorem]
\label{rem:chow-iv-relative-meeks-yau}
\source{chow-et-al-ricci-flow-part-iv:page-0260}
\dcref{ch33:33.26}
\group{Minimal disks}

The same least-area conclusion holds when the allowed boundary is any selected
union of boundary components, provided the corresponding disk class is
nonempty.
\end{remark}

\begin{lemma}[First variation bound for a minimizing disk]
\label{lem:chow-iv-minimizing-disk-first-variation}
\source{chow-et-al-ricci-flow-part-iv:page-0261,chow-et-al-ricci-flow-part-iv:page-0262,chow-et-al-ricci-flow-part-iv:page-0263,chow-et-al-ricci-flow-part-iv:page-0264}
\uses{thm:chow-iv-meeks-yau}
\dcref{ch33:33.27}
\group{Minimal disks}


Let \(D(t_0)\) be a least-area disk normal to the moving boundary torus, let
\(k\) be the geodesic curvature of its boundary in the disk, and let \(V^0\)
be the normal velocity of that boundary.  For all sufficiently large \(t_0\),
\[
 \frac{d^+A_{\mathfrak T_A}}{dt}(t_0)\le -2\pi
 +\int_{\partial D(t_0)}(k+V^0)\,ds
 +\int_{D(t_0)}\left(\frac23r-\frac12R\right)\,dA.
\]
\end{lemma}

\begin{proof}
\sketch
Use the disk at \(t_0\) to construct smoothly varying comparison disks ending
on the moving tori.  The normalized Ricci-flow variation of their area is the
interior metric variation plus the boundary velocity term.  Expressing the
tangential Ricci trace through scalar and sectional curvature, the Gauss
equation and minimality give \(K\le\operatorname{sec}(TD)\).  Gauss--Bonnet on
the disk contributes \(-2\pi+\int_{\partial D}k\,ds\), yielding the inequality.
\end{proof}

\begin{lemma}[Evolution of the area of a minimizing disk]
\label{lem:chow-iv-minimizing-disk-evolution}
\source{chow-et-al-ricci-flow-part-iv:page-0264,chow-et-al-ricci-flow-part-iv:page-0265}
\uses{lem:chow-iv-minimizing-disk-first-variation}
\dcref{ch33:33.28}
\group{Minimal disks}


For every \(A\in(0,\widetilde A)\) and \(\varepsilon>0\), at all sufficiently
large times a minimizing disk \(D(t)\) satisfies
\[
 \frac{d^+A_{\mathfrak T_A}}{dt}(t)\le-2\pi
 +(1+\varepsilon)^{1/2}L_{g(t)}(\partial D(t))
 -(1+\varepsilon)^{-1/2}A_{\mathfrak T_A}(t).
\]
\end{lemma}

\begin{proof}
\sketch
The average and minimum scalar curvatures tend to \(-6\), so the interior term
in the first-variation bound is at most
\(-(1-\varepsilon)A_{\mathfrak T_A}\).  In an exact cusp the relevant boundary
geodesic curvature is one; smooth convergence of the piece to a cusp makes its
integral at most \((1+\varepsilon)L\).  The normal velocity tends uniformly to
zero.  Renaming \(\varepsilon\) gives the displayed estimate.
\end{proof}

\begin{lemma}[Embedded disks meet almost every torus slice transversely]
\label{lem:chow-iv-disk-torus-slice-transversality}
\source{chow-et-al-ricci-flow-part-iv:page-0265,chow-et-al-ricci-flow-part-iv:page-0266}
\dcref{ch33:33.29}
\group{Cusp comparison}

If a smooth embedded disk meets a torus boundary in a product collar
\(T^2\times[a,b)\), then for almost every \(z\in(a,b)\) it is transverse to
\(T^2\times\{z\}\), and the intersection is a finite disjoint union of smooth
embedded loops.
\end{lemma}

\begin{proof}
\sketch
On the inverse image of the collar, compose the disk embedding with projection
to \([a,b)\).  Sard's theorem makes almost every \(z\) a regular value.  Its
preimage is then a compact one-manifold without boundary, hence a finite union
of circles, and embeddedness gives the assertion for the image.
\end{proof}

\begin{lemma}[Area decay when boundary length is controlled]
\label{lem:chow-iv-area-decay-short-boundary}
\source{chow-et-al-ricci-flow-part-iv:page-0266}
\uses{lem:chow-iv-minimizing-disk-evolution}
\dcref{ch33:33.30}
\group{Minimal disks}


If, at a sufficiently large time,
\(L_{g(t)}(\partial D(t))\le(1+\varepsilon)^{-1}A_{\mathfrak T_A}(t)\), then
\[
 \frac{d^+A_{\mathfrak T_A}}{dt}(t)\le-2\pi.
\]
\end{lemma}

\begin{definition}[Almost hyperbolic cusp region]
\label{def:chow-iv-almost-hyperbolic-cusp-region}
\source{chow-et-al-ricci-flow-part-iv:page-0266,chow-et-al-ricci-flow-part-iv:page-0267}
\dcref{ch33:33.31}
\group{Cusp comparison}

For \(0<B<A<\widetilde A\), the \emph{almost hyperbolic cusp region}
\(\mathcal C_{A,B}(t)\) is the region bounded by a chosen component
\(T_A(t)\subset\partial\mathcal M_A(t)\) and the corresponding component of
\(\partial\mathcal M_B(t)\).  Its parallel tori are denoted
\(T_{A,\rho}(t)=\{d_{g(t)}(\,cdot\,,T_A(t))=\rho\}\).
\end{definition}

\begin{definition}[Depth of an almost cusp]
\label{def:chow-iv-almost-cusp-depth}
\source{chow-et-al-ricci-flow-part-iv:page-0267}
\uses{def:chow-iv-almost-hyperbolic-cusp-region}
\dcref{ch33:33.32}
\group{Cusp comparison}


Define \(\rho_{A,B}(t)\) to be the largest \(\sigma\) such that
\(T_{A,\rho}(t)\subset\mathcal C_{A,B}(t)\) for every
\(\rho\in[0,\sigma]\).
\end{definition}

\begin{remark}[Almost cusps have arbitrary depth]
\label{rem:chow-iv-almost-cusps-arbitrary-depth}
\source{chow-et-al-ricci-flow-part-iv:page-0267}
\uses{def:chow-iv-almost-cusp-depth}
\dcref{ch33:33.33}
\group{Cusp comparison}


For fixed \(A\) and finite \(N\), choosing \(B<A\) sufficiently small and then
\(t\) sufficiently large ensures \(\rho_{A,B}(t)\ge N\).
\end{remark}

\begin{lemma}[Monotonicity-type formula for length]
\label{lem:chow-iv-cusp-length-monotonicity}
\source{chow-et-al-ricci-flow-part-iv:page-0267,chow-et-al-ricci-flow-part-iv:page-0268,chow-et-al-ricci-flow-part-iv:page-0269,chow-et-al-ricci-flow-part-iv:page-0270}
\uses{lem:chow-iv-disk-torus-slice-transversality, def:chow-iv-almost-cusp-depth}
\dcref{ch33:33.34}
\group{Cusp comparison}


Let \(L(\rho)\) be the total length of the intersection of a minimizing disk
with \(T_{A,\rho}\).  For every sufficiently small \(\varepsilon>0\), at all
sufficiently large times and for \(0<\bar\rho\le\rho_{A,B}(t)\),
\[
 \frac d{d\bar\rho}\left(
 \frac{e^{\varepsilon\bar\rho}}{1-e^{-\bar\rho}}
 \int_0^{\bar\rho}L(\rho)\,d\rho\right)\ge0.
\]
\end{lemma}

\begin{proof}
\sketch
The coarea formula bounds the area swept out by the disk below by
\(\int L\).  Replace an innermost essential intersection loop by the normal
annulus over it; minimality bounds the disk area by this comparison annulus.
In an almost cusp, lengths along normal rays differ from exact exponential
decay by an arbitrarily small error.  Integrating that error bounds
\(\int_0^{\bar\rho}L\) in terms of \(L(\bar\rho)\); differentiating the displayed
quantity then gives a nonnegative derivative.
\end{proof}

\begin{corollary}[Length bound by swept area]
\label{cor:chow-iv-cusp-length-by-swept-area}
\source{chow-et-al-ricci-flow-part-iv:page-0271}
\uses{lem:chow-iv-cusp-length-monotonicity}
\dcref{ch33:33.35}
\group{Cusp comparison}


At sufficiently large times,
\[
 L(0)\le
 \frac{e^{\varepsilon\bar\rho}}{1-e^{-\bar\rho}}
 \int_0^{\bar\rho}L(\rho)\,d\rho
 \le\frac{e^{\varepsilon\bar\rho}}{1-e^{-\bar\rho}}A(\bar\rho)
\]
for every \(0<\bar\rho<\rho_{A,B}(t)\), where \(A(\bar\rho)\) is the disk area
between depths zero and \(\bar\rho\).
\end{corollary}

\begin{corollary}[Area almost bounds boundary length]
\label{cor:chow-iv-area-almost-bounds-length}
\source{chow-et-al-ricci-flow-part-iv:page-0271}
\uses{cor:chow-iv-cusp-length-by-swept-area, rem:chow-iv-almost-cusps-arbitrary-depth}
\dcref{ch33:33.36}
\group{Cusp comparison}


For every \(A\in(0,\widetilde A)\) and \(\delta>0\), one may choose \(B<A\)
and a time \(T_\delta\) such that for \(t\ge T_\delta\),
\[
 L_{g(t)}(\partial D(t))\le(1+\delta)A(N)
 \le(1+\delta)A(\rho_{A,B}(t)),
 \quad N=-\log\left(1-\frac1{\sqrt{1+\delta}}\right).
\]
\end{corollary}

\begin{lemma}[A short intersection loop deep in the cusp]
\label{lem:chow-iv-short-loop-deep-cusp}
\source{chow-et-al-ricci-flow-part-iv:page-0271,chow-et-al-ricci-flow-part-iv:page-0272}
\uses{cor:chow-iv-area-almost-bounds-length}
\dcref{ch33:33.37}
\group{Cusp comparison}


Assume the boundary length is larger than
\((1+\varepsilon)^{-1}A_{\mathfrak T_A}(t_\diamond)\).  There is a constant
\(\Lambda(\varepsilon)\) such that, in every sufficiently long interval
\([\rho^\#,\rho^*]\) of cusp depths, some regular depth \(\rho_\diamond\) satisfies
\[
 L(\rho_\diamond)\le(1+2\varepsilon)
 e^{-(\rho^*-\rho_\diamond)}L(0).
\]
\end{lemma}

\begin{proof}
\sketch
Choose a regular depth minimizing \(e^{-\rho}L(\rho)\) on the interval.  Coarea,
the assumed lower bound for \(L(0)\), and the preceding area estimate give
\((e^{\rho^*}-e^{\rho^\#})e^{-\rho_\diamond}L(\rho_\diamond)
\le(1+\varepsilon)L(0)\).  Taking
\(\Lambda\ge\log(\varepsilon^{-1}+2)\) gives the claimed constant.
\end{proof}

\begin{remark}[Depth required for the short-loop lemma]
\label{rem:chow-iv-short-loop-depth-requirement}
\source{chow-et-al-ricci-flow-part-iv:page-0271}
\uses{lem:chow-iv-short-loop-deep-cusp, rem:chow-iv-almost-cusps-arbitrary-depth}
\dcref{ch33:33.38}
\group{Cusp comparison}


The requirement \(\rho^\#+\Lambda\le\rho_{A,B}(t_\diamond)\) is achieved by
choosing the inner truncation area \(B\) sufficiently small.
\end{remark}

\begin{lemma}[Annulus estimate in a flat torus]
\label{lem:chow-iv-flat-torus-annulus-area}
\source{chow-et-al-ricci-flow-part-iv:page-0273,chow-et-al-ricci-flow-part-iv:page-0274}
\dcref{ch33:33.39}
\group{Cusp comparison}

Let \(T^2\) be flat.  If an embedded geodesic loop \(L\) and a smooth embedded
loop \(S\) represent the same fundamental-group element, then a geodesic loop
\(L'\) parallel to \(L\) and an immersed annulus \(\mathcal A\) bounded by
\(S\cup L'\) exist with
\[
 \operatorname{Area}(\mathcal A)
 \le\operatorname{Length}(L)\operatorname{Length}(S).
\]
\end{lemma}

\begin{proof}
\sketch
Quotient the universal cover \(\mathbb R^2\) by the deck translation
represented by \(L\).  This is a flat cylinder of circumference
\(\operatorname{Length}(L)\).  The lift of \(S\) lies in a band of width at most
\(\operatorname{Length}(S)\).  A meridian in that band and the lifted loop
bound an immersed annulus of area at most the area of the band.  Projecting to
\(T^2\) proves the estimate.
\end{proof}

\begin{proposition}[Length is small]
\label{prop:chow-iv-minimizing-disk-boundary-small}
\source{chow-et-al-ricci-flow-part-iv:page-0272,chow-et-al-ricci-flow-part-iv:page-0273,chow-et-al-ricci-flow-part-iv:page-0274,chow-et-al-ricci-flow-part-iv:page-0275}
\uses{lem:chow-iv-short-loop-deep-cusp, lem:chow-iv-flat-torus-annulus-area, cor:chow-iv-area-almost-bounds-length}
\dcref{ch33:33.40}
\group{Cusp comparison}


For every \(\varepsilon_1>0\), the outer truncation area \(A\) may be chosen
small enough that, at all sufficiently large times for which the long-boundary
alternative holds, every area-minimizing disk satisfies
\[
 L_{g(t)}(\partial D(t))\le\varepsilon_1.
\]
\end{proposition}

\begin{proof}
\sketch
At the depth supplied by the short-loop lemma, replace an essential
intersection loop by a shortest geodesic in its flat torus.  Join the two by
the flat-torus annulus and extend the geodesic back along normal cusp rays to
obtain a comparison disk.  Almost-hyperbolic cusp estimates bound the two new
annuli by a constant times the length \(L_0\) of the corresponding geodesic on
the area-\(A\) model torus.  Minimality and the area-almost-bounds-length
corollary then bound \(L(0)\) by a multiple of \(L_0\), plus a term that can be
absorbed.  Since \(L_0\to0\) with \(A\), the result follows.
\end{proof}

\begin{lemma}[Area decay in the long-boundary case]
\label{lem:chow-iv-area-decay-long-boundary}
\source{chow-et-al-ricci-flow-part-iv:page-0275}
\uses{lem:chow-iv-minimizing-disk-evolution, prop:chow-iv-minimizing-disk-boundary-small}
\dcref{ch33:33.41}
\group{Minimal disks}


For sufficiently small \(\varepsilon>0\), sufficiently small truncation area,
and all sufficiently large times at which
\(L(\partial D)>(1+\varepsilon)^{-1}A_{\mathfrak T_A}\),
\[
 \frac{d^+A_{\mathfrak T_A}}{dt}\le-2\pi+\varepsilon.
\]
\end{lemma}

\begin{proof}
\sketch
Choose the boundary-length bound in Proposition~\ref{prop:chow-iv-minimizing-disk-boundary-small}
to be \(\varepsilon/3\), and substitute it into the minimizing-disk evolution
inequality.  Its negative area term may be discarded; after decreasing
\(\varepsilon\) initially, the remaining boundary contribution is at most
\(\varepsilon\).
\end{proof}

\begin{claim}[Exercise 33.20: collapse improves to zero]
\label{ex:chow-iv-collapsed-complement-injectivity-zero}
\dcref{ch33:33.20}
\group{Stability of hyperbolic limits}
\source{chow-et-al-ricci-flow-part-iv:page-0258}
\uses{prop:chow-iv-hyperbolic-collapsed-decomposition}
For the finite family of stable hyperbolic pieces in the decomposition,
\[
 \max_{x\in\mathcal M\setminus\bigcup_{a=0}^N\mathcal M_{a,A_a(t)}(t)}
 \operatorname{inj}_{g(t)}(x)\longrightarrow0
 \quad\text{as }t\to\infty.
\]
\end{claim}
